\chapter{Case Study and Threat Modeling}\label{ch:threatModeling}

As a first step, two real-world attacks that occurred within the \acp{NPM}
Software Supply Chain (\acp{SSC}) were selected and analyzed.

The study of these cases made it possible to gain a deeper understanding
of the main characteristics of these attacks,
providing the basis for the subsequent definition 
of the metrics necessary to identify these types of threats,
in the \nameref{ch:metrics} (\Cref{ch:metrics}).

\section{Crypto related attack - September 2025}

The first attack, which occurred on September 8 2025,
compromised 18 highly popular \acp{NPM} packages, including \texttt{chalk},
\texttt{debug}, and \texttt{ansi-colors} \cite{henig2025breakdownwidespreadnpmsupplychainattack}.

The attack was made possible by compromising the account of a maintainer of one of the packages via a phishing email.

After obtaining the account credentials,
the attacker published malicious versions of the affected packages,
which contained malware designed to steal cryptocurrencies
by intercepting and manipulating users’ browser transactions.

The malware's operation can be divided into these main phases \cite{matthee2025npmsupplychainattack}.
First, the malicious code checked for the presence of a Web3 wallet interface
in the browser (e.g., \texttt{window.ethereum}).

If present, it hooked one of these specific functions
(\texttt{eth\_send\-Transaction}, \texttt{solana\_sign\-Transaction},
\texttt{solana\_sign\-And\-Send\-Transaction})
with the goal of replacing the transaction's destination address
with one controlled by the attacker; 
this is done by intercepting network calls (using the \texttt{fetch} or \texttt{XMLHttpRequest} \acp{APIs}).

The compromised package versions remained online
for approximately two hours before being detected by developers and removed by \acp{NPM}.
As a result, the overall impact of the attack was limited.

\section{Information theft attack - November 2025}

The second attack considered in this study occurred on November 24 2025.
Within a few hours, over 1,000 \acp{NPM} packages were infected,
all using the same technique; again relying on compromised maintainer accounts
\cite{ravie2025secondsha1hulud}.

The attack, known as \textit{Sha1-Hulud}, introduced the Bun runtime
as a dependency in the packages,
adding a \texttt{preinstall} script (\texttt{node setup\_bun.js})
along with a highly obfuscated file named \texttt{bun\_environment.js},
with a size exceeding 10 MB \cite{helixguard2025malicioussha1hulud}.

When executed, the script downloaded and ran TruffleHog
to scan the local system,
stealing sensitive information such as \acp{NPM} tokens,
\acp{AWS}, \acp{GCP}, and Azure credentials, and environment variables.

